% Source: https://artofproblemsolving.com/community/c3h1696154_graphing_inequalities
% Stable identity: aops-topic-1696154-q01
Solve each inequality for $y$:
\[
2x-3y+7\le0\iff y\ge\dfrac{2x+7}{3},\qquad x+3y\ge-1\iff y\ge-\dfrac{x+1}{3}.
\]
So the region consists of the points on or above \emph{both} lines $y=\dfrac{2x+7}{3}$ (slope $\dfrac{2}{3}$, through $(-2,1)$ and $(1,3)$) and $y=-\dfrac{x+1}{3}$ (slope $-\dfrac{1}{3}$, through $(-1,0)$ and $(2,-1)$), boundaries included (both inequalities are non-strict). The lines meet where $2x-3y=-7$ and $x+3y=-1$; adding gives $3x=-8$, so $x=-\dfrac{8}{3}$ and $y=\dfrac{5}{9}$.

The region is the $\boxed{\text{wedge above both lines with vertex }V=\left(-\dfrac{8}{3},\dfrac{5}{9}\right)}$: to the left of the vertex its lower edge is the line $x+3y=-1$, and to the right it is the line $2x-3y+7=0$. For example $(0,3)$ is in it ($-9+7\le0$ and $9\ge-1$), while the origin is not ($7>0$).

\begin{asy}
usepackage("amsmath");
size(8cm);
defaultpen(fontsize(10pt));
import graph;
real f1(real x){return (2x+7)/3;} real f2(real x){return -(x+1)/3;}
real xa=-6,xb=4,ya=-3,yb=6;
fill((xa,f2(xa))--(-8/3,5/9)--(xb,f1(xb))--(xb,yb)--(xa,yb)--cycle,lightblue+opacity(0.6));
draw((xa,0)--(xb+0.4,0),gray,Arrow(TeXHead));draw((0,ya)--(0,yb+0.4),gray,Arrow(TeXHead));
draw((-5,f1(-5))--(xb,f1(xb)),blue+linewidth(0.9));
draw((xa,f2(xa))--(xb,f2(xb)),red+linewidth(0.9));
dot((-8/3,5/9));
label("$x$",(xb+0.4,0),E);label("$y$",(0,yb+0.4),N);
label("$2x-3y+7=0$",(2.2,f1(2.2)),SE,blue);
label("$x+3y=-1$",(3.5,f2(3.5)),S,red);
label("$V$",(-8/3,5/9),dir(97));
for(int k=-6;k<=4;k+=2) if(k!=0 && k!=-4) {draw((k,-0.12)--(k,0.12),gray);label("$"+string(k)+"$",(k,-0.12),S,fontsize(8pt));}
for(int k=-2;k<=4;k+=2) if(k!=0) {draw((-0.12,k)--(0.12,k),gray);label("$"+string(k)+"$",(0.12,k),E,fontsize(8pt));}
label("\mbox{}",(min(currentpicture,user=true)+max(currentpicture,user=true))/2,black); // keeps the text after the figure black when LaTeX typesets the labels inline
\end{asy}
